% -------------------------------------------------
% Anexo H. Endpoints de la API REST
% -------------------------------------------------
\section*{Anexo H. Especificación detallada de endpoints de la API REST}
\addcontentsline{toc}{section}{Anexo H. Especificación detallada de endpoints de la API REST}
\label{anexo:endpoints}

Esta sección presenta la descripción detallada de los endpoints expuestos por la API REST
del sistema, organizados por dominios funcionales. Para cada grupo se especifican el método
HTTP, la ruta, la acción asociada, una breve descripción, el tipo de autenticación requerida
y el esquema principal de respuesta.

\subsection*{Autenticación y usuarios}

\begin{table}[H]
\centering
\footnotesize
\setlength{\tabcolsep}{4pt}
\renewcommand{\arraystretch}{1.25}
\begin{tabular}{|p{1.4cm}|p{4.1cm}|p{3.2cm}|p{5.8cm}|p{1cm}|}
\hline
\textbf{Método} & \textbf{Ruta} & \textbf{Recurso / acción} & \textbf{Descripción} & \textbf{Auth} \\
\hline

POST & \url{/api/v1/auth/login} & Inicio de sesión &
Autentica al usuario y devuelve un token JWT. &
No \\
\hline

POST & \url{/api/v1/auth/register} & Registro de usuario &
Crea una nueva cuenta de usuario. &
No \\
\hline

PATCH & \url{/api/v1/users/me} & Actualizar perfil &
Actualiza datos del usuario autenticado. &
Sí \\
\hline

DELETE & \url{/api/v1/users/me} & Eliminar cuenta &
Desactiva la cuenta del usuario autenticado. &
Sí \\
\hline

GET & \url{/api/v1/users/me/learning-goals} & Mis metas de aprendizaje &
Obtiene las metas de aprendizaje del usuario con paginación. &
Sí \\
\hline

GET & \url{/api/v1/users/me/pomodoro-preferences} & Obtener preferencias Pomodoro &
Obtiene la configuración de preferencias Pomodoro del usuario. &
Sí \\
\hline

PUT & \url{/api/v1/users/me/pomodoro-preferences} &
Crear/actualizar preferencias Pomodoro &
Crea o actualiza las preferencias Pomodoro del usuario. &
Sí \\
\hline

\end{tabular}
\caption{Endpoints del dominio de autenticación y usuarios.}
\end{table}

% -------------------------------------------------
\subsection*{Módulo Aprender a Aprender}

\subsubsection*{Metas de aprendizaje (\textit{learning\_goals})}

\begin{table}[H]
\centering
\footnotesize
\setlength{\tabcolsep}{4pt}
\renewcommand{\arraystretch}{1.25}
\begin{tabular}{|p{1.4cm}|p{4.1cm}|p{3.2cm}|p{5.8cm}|p{1cm}|}
\hline
\textbf{Método} & \textbf{Ruta} & \textbf{Recurso / acción} & \textbf{Descripción} & \textbf{Auth} \\
\hline

POST & \url{/api/v1/learning-goals} & Crear meta de aprendizaje &
Crea una nueva meta de aprendizaje. &
Sí \\
\hline

GET & \url{/api/v1/learning-goals/{id}} & Obtener meta por ID &
Obtiene los detalles de una meta de aprendizaje. &
Sí \\
\hline

GET & \url{/api/v1/learning-goals/{id}/objectives} & Objetivos de una meta &
Lista los objetivos de una meta con filtros y paginación. &
Sí \\
\hline

PATCH & \url{/api/v1/learning-goals/{id}} & Actualizar meta &
Actualiza los datos de una meta de aprendizaje. &
Sí \\
\hline

DELETE & \url{/api/v1/learning-goals/{id}} & Eliminar meta &
Elimina una meta de aprendizaje. &
Sí \\
\hline

POST & \url{/api/v1/learning-goals/{id}/convert-to-roadmap} & Convertir a roadmap &
Convierte una meta de aprendizaje en un roadmap. &
Sí \\
\hline

\end{tabular}
\caption{Endpoints para la gestión de metas de aprendizaje.}
\end{table}

% -------------------------------------------------
\subsubsection*{Objetivos (\textit{objectives})}

\begin{table}[H]
\centering
\footnotesize
\setlength{\tabcolsep}{4pt}
\renewcommand{\arraystretch}{1.25}
\begin{tabular}{|p{1.4cm}|p{4.1cm}|p{3.2cm}|p{5.8cm}|p{1cm}|}
\hline
\textbf{Método} & \textbf{Ruta} & \textbf{Recurso / acción} & \textbf{Descripción} & \textbf{Auth} \\
\hline

POST & \url{/api/v1/objectives} & Crear objetivo &
Crea un nuevo objetivo de aprendizaje. &
Sí \\
\hline

GET & \url{/api/v1/objectives/{id}} & Obtener objetivo &
Obtiene los detalles de un objetivo. &
Sí \\
\hline

GET & \url{/api/v1/objectives/{id}/tasks} & Tareas de un objetivo &
Lista tareas de un objetivo con filtros y paginación. &
Sí \\
\hline

PATCH & \url{/api/v1/objectives/{id}} & Actualizar objetivo &
Actualiza un objetivo existente. &
Sí \\
\hline

DELETE & \url{/api/v1/objectives/{id}} & Eliminar objetivo &
Elimina un objetivo. &
Sí \\
\hline

GET & \url{/api/v1/objectives/{id}/kanban} & Obtener tablero Kanban &
Obtiene el tablero Kanban de un objetivo. &
Sí \\
\hline

GET & \url{/api/v1/objectives/{id}/kanban/column} & Columna Kanban &
Obtiene tareas paginadas de una columna específica. &
Sí \\
\hline

PATCH & \url{/api/v1/objectives/{id}/kanban/move} & Mover tarea en Kanban &
Mueve una tarea entre columnas del tablero Kanban. &
Sí \\
\hline

\end{tabular}
\caption{Endpoints para la gestión de objetivos.}
\end{table}

% -------------------------------------------------
\subsubsection*{Tareas (\textit{tasks})}

\begin{table}[H]
\centering
\footnotesize
\setlength{\tabcolsep}{4pt}
\renewcommand{\arraystretch}{1.25}
\begin{tabular}{|p{1.4cm}|p{4.1cm}|p{3.2cm}|p{5.8cm}|p{1cm}|}
\hline
\textbf{Método} & \textbf{Ruta} & \textbf{Recurso / acción} & \textbf{Descripción} & \textbf{Auth} \\
\hline

POST & \url{/api/v1/tasks} & Crear tarea &
Crea una nueva tarea asociada a un objetivo. &
Sí \\
\hline

GET & \url{/api/v1/tasks/{id}} & Obtener tarea &
Obtiene los detalles de una tarea. &
Sí \\
\hline

PATCH & \url{/api/v1/tasks/{id}} & Actualizar tarea &
Actualiza los campos de una tarea. &
Sí \\
\hline

PATCH & \url{/api/v1/tasks/{id}/status} & Actualizar estado &
Modifica el estado de una tarea. &
Sí \\
\hline

DELETE & \url{/api/v1/tasks/{id}} & Eliminar tarea &
Elimina una tarea. &
Sí \\
\hline

\end{tabular}
\caption{Endpoints para la gestión de tareas.}
\end{table}

% -------------------------------------------------
\subsubsection*{Notas de tarea (\textit{task\_notes})}

\begin{table}[H]
\centering
\footnotesize
\setlength{\tabcolsep}{4pt}
\renewcommand{\arraystretch}{1.25}
\begin{tabular}{|p{1.4cm}|p{4.1cm}|p{3.2cm}|p{5.8cm}|p{1cm}|}
\hline
\textbf{Método} & \textbf{Ruta} & \textbf{Recurso / acción} & \textbf{Descripción} & \textbf{Auth} \\
\hline

POST & \url{/api/v1/tasks/{task\_id}/notes} & Agregar nota &
Crea una nota asociada a una tarea. &
Sí \\
\hline

GET & \url{/api/v1/tasks/{task\_id}/notes} & Listar notas &
Lista todas las notas de una tarea. &
Sí \\
\hline

PATCH & \url{/api/v1/tasks/{task\_id}/notes/{note\_id}} & Actualizar nota &
Actualiza una nota de tarea. &
Sí \\
\hline

DELETE & \url{/api/v1/tasks/{task\_id}/notes/{note\_id}} & Eliminar nota &
Elimina una nota de tarea. &
Sí \\
\hline

\end{tabular}
\caption{Endpoints para la gestión de notas de tarea.}
\end{table}

% -------------------------------------------------
\subsubsection*{Recursos de tarea (\textit{task\_resources})}

\begin{table}[H]
\centering
\footnotesize
\setlength{\tabcolsep}{4pt}
\renewcommand{\arraystretch}{1.25}
\begin{tabular}{|p{1.4cm}|p{4.1cm}|p{3.2cm}|p{5.8cm}|p{1cm}|}
\hline
\textbf{Método} & \textbf{Ruta} & \textbf{Recurso / acción} & \textbf{Descripción} & \textbf{Auth} \\
\hline

POST & \url{/api/v1/tasks/{task\_id}/resources} & Agregar recurso &
Crea un recurso asociado a una tarea. &
Sí \\
\hline

GET & \url{/api/v1/tasks/{task\_id}/resources} & Listar recursos &
Lista todos los recursos de una tarea. &
Sí \\
\hline

PATCH & \url{/api/v1/tasks/{task\_id}/resources/{resource\_id}} & Actualizar recurso &
Actualiza un recurso de tarea. &
Sí \\
\hline

DELETE & \url{/api/v1/tasks/{task\_id}/resources/{resource\_id}} & Eliminar recurso &
Elimina un recurso de tarea. &
Sí \\
\hline

\end{tabular}
\caption{Endpoints para la gestión de recursos asociados a tareas.}
\end{table}

% -------------------------------------------------
\subsubsection*{Autoevaluaciones (\textit{self\_evaluations})}

\begin{table}[H]
\centering
\footnotesize
\setlength{\tabcolsep}{4pt}
\renewcommand{\arraystretch}{1.25}
\begin{tabular}{|p{1.4cm}|p{4.1cm}|p{3.2cm}|p{5.8cm}|p{1cm}|}
\hline
\textbf{Método} & \textbf{Ruta} & \textbf{Recurso / acción} & \textbf{Descripción} & \textbf{Auth} \\
\hline

POST & \url{/api/v1/self-evaluations} & Crear autoevaluación &
Registra una autoevaluación para una tarea completada. &
Sí \\
\hline

GET & \url{/api/v1/self-evaluations/{evaluation\_id}} & Obtener autoevaluación &
Obtiene una autoevaluación por identificador. &
Sí \\
\hline

GET & \url{/api/v1/self-evaluations/by-user/{user\_id}} & Autoevaluaciones por usuario &
Lista autoevaluaciones de un usuario con filtros y paginación. &
Sí \\
\hline

\end{tabular}
\caption{Endpoints para la gestión de autoevaluaciones.}
\end{table}

% -------------------------------------------------
\subsection*{Módulo Escucha Activa}

\subsubsection*{Sesiones de juego (\textit{game\_sessions})}

\begin{table}[H]
\centering
\footnotesize
\setlength{\tabcolsep}{4pt}
\renewcommand{\arraystretch}{1.25}
\begin{tabular}{|p{1.4cm}|p{4.1cm}|p{3.2cm}|p{5.8cm}|p{1cm}|}
\hline
\textbf{Método} & \textbf{Ruta} & \textbf{Recurso / acción} &
\textbf{Descripción} & \textbf{Auth} \\
\hline

GET & \url{/api/v1/game-sessions} & Listar sesiones &
Lista sesiones de juego del usuario con paginación. &
Sí \\
\hline

POST & \url{/api/v1/game-sessions} & Crear sesión de juego &
Crea una nueva sesión de juego de escucha. &
Sí \\
\hline

POST & \url{/api/v1/game-sessions/{session\_id}/start} & Iniciar sesión &
Inicia una sesión y devuelve la primera ronda. &
Sí \\
\hline

GET & \url{/api/v1/game-sessions/{session\_id}} & Obtener detalles &
Obtiene el detalle de una sesión, incluida su configuración. &
Sí \\
\hline

PATCH & \url{/api/v1/game-sessions/{session\_id}} & Actualizar sesión &
Actualiza nombre y/o estado de la sesión. &
Sí \\
\hline

PATCH & \url{/api/v1/game-sessions/{session\_id}/config} & Actualizar configuración &
Actualiza parcialmente la configuración de la sesión. &
Sí \\
\hline

DELETE & \url{/api/v1/game-sessions/{session\_id}} & Eliminar sesión &
Elimina una sesión de juego. &
Sí \\
\hline

GET & \url{/api/v1/game-sessions/{session\_id}/rounds/current} & Ronda actual &
Obtiene la ronda actual de la sesión. &
Sí \\
\hline

GET & \url{/api/v1/game-sessions/{session\_id}/rounds/{round\_number}} & Ronda específica &
Obtiene una ronda específica por número. &
Sí \\
\hline

POST & \url{/api/v1/game-sessions/{session\_id}/rounds/next} & Siguiente ronda &
Avanza a la siguiente ronda de la sesión. &
Sí \\
\hline

POST & \url{/api/v1/game-sessions/{session\_id}/rounds/{round\_number}/attempt} &
Enviar intento &
Envía un intento de respuesta para una ronda. &
Sí \\
\hline

POST & \url{/api/v1/game-sessions/{session\_id}/rounds/{round\_number}/audio/replay} &
Repetición de audio &
Incrementa el contador de reproducción de audio si está bajo el límite. &
Sí \\
\hline

POST & \url{/api/v1/game-sessions/{session\_id}/finish} & Finalizar sesión &
Marca la sesión como finalizada. &
Sí \\
\hline

GET & \url{/api/v1/game-sessions/{session\_id}/result} & Resultado de sesión &
Obtiene el resultado de una sesión terminada. &
Sí \\
\hline

\end{tabular}
\caption{Endpoints para la gestión de sesiones de juego del módulo Escucha Activa.}
\end{table}

% -------------------------------------------------
\subsubsection*{Desafíos de escucha (\textit{challenges})}

\begin{table}[H]
\centering
\footnotesize
\setlength{\tabcolsep}{4pt}
\renewcommand{\arraystretch}{1.25}
\begin{tabular}{|p{1.4cm}|p{4.1cm}|p{3.2cm}|p{5.8cm}|p{1cm}|}
\hline
\textbf{Método} & \textbf{Ruta} & \textbf{Recurso / acción} &
\textbf{Descripción} & \textbf{Auth} \\
\hline

POST & \url{/api/v1/challenges/generate} & Generar desafío &
Genera un nuevo desafío de escucha usando un modelo LLM. &
No \\
\hline

GET & \url{/api/v1/challenges/{challenge\_id}/audio} & Audio de desafío &
Obtiene o crea el audio asociado a un desafío y devuelve su URL. &
No \\
\hline

\end{tabular}
\caption{Endpoints para la gestión de desafíos de escucha.}
\end{table}

% -------------------------------------------------
\subsection*{Roadmaps (planes de aprendizaje)}

\begin{table}[H]
\centering
\footnotesize
\setlength{\tabcolsep}{4pt}
\renewcommand{\arraystretch}{1.25}
\begin{tabular}{|p{1.4cm}|p{4.1cm}|p{3.2cm}|p{5.8cm}|p{1cm}|}
\hline
\textbf{Método} & \textbf{Ruta} & \textbf{Recurso / acción} &
\textbf{Descripción} & \textbf{Auth} \\
\hline

POST & \url{/api/v1/roadmap} & Crear roadmap &
Crea un nuevo plan de aprendizaje. &
Sí \\
\hline

GET & \url{/api/v1/roadmap/mine} & Mis roadmaps &
Obtiene los planes de aprendizaje del usuario autenticado con paginación. &
Sí \\
\hline

GET & \url{/api/v1/roadmap/public} & Roadmaps públicos &
Lista planes públicos con filtros por título, autor, fechas y número de pasos. &
Sí \\
\hline

GET & \url{/api/v1/roadmap/{id}} & Obtener roadmap &
Obtiene los detalles de un plan de aprendizaje. &
Sí \\
\hline

PATCH & \url{/api/v1/roadmap/{id}} & Actualizar roadmap &
Actualiza la información de un plan de aprendizaje. &
Sí \\
\hline

POST & \url{/api/v1/roadmap/{id}/copy} & Copiar roadmap &
Realiza una copia (clonación) de un plan de aprendizaje. &
Sí \\
\hline

POST & \url{/api/v1/roadmap/{id}/convert-to-learning-goal} & Convertir a meta &
Convierte un roadmap en una meta de aprendizaje accionable. &
Sí \\
\hline

DELETE & \url{/api/v1/roadmap/{id}} & Eliminar roadmap &
Elimina un plan de aprendizaje. &
Sí \\
\hline

\end{tabular}
\caption{Endpoints para la gestión de roadmaps.}
\end{table}

% -------------------------------------------------
\subsection*{Módulos del sistema}

\begin{table}[H]
\centering
\footnotesize
\setlength{\tabcolsep}{4pt}
\renewcommand{\arraystretch}{1.25}
\begin{tabular}{|p{1.4cm}|p{4.1cm}|p{3.2cm}|p{5.8cm}|p{1cm}|}
\hline
\textbf{Método} & \textbf{Ruta} & \textbf{Recurso / acción} &
\textbf{Descripción} & \textbf{Auth} \\
\hline

POST & \url{/api/v1/modules} & Crear módulo &
Crea un nuevo módulo del sistema (requiere rol administrador). &
Sí (admin) \\
\hline

GET & \url{/api/v1/modules} & Listar módulos &
Lista los módulos disponibles con paginación. &
Sí \\
\hline

GET & \url{/api/v1/modules/{id}} & Obtener módulo &
Obtiene los detalles de un módulo. &
Sí \\
\hline

PATCH & \url{/api/v1/modules/{id}} & Actualizar módulo &
Actualiza la información de un módulo (requiere rol administrador). &
Sí (admin) \\
\hline

DELETE & \url{/api/v1/modules/{id}} & Eliminar módulo &
Elimina un módulo del sistema (requiere rol administrador). &
Sí (admin) \\
\hline

\end{tabular}
\caption{Endpoints para la gestión de módulos del sistema.}
\end{table}